% \documentclass[12pt]{article}
\documentclass[12pt,fleqn]{article}
\usepackage[utf8]{inputenc}
\usepackage[russian]{babel}
\usepackage{graphicx}
\usepackage{array}
\usepackage[a4paper, left=1cm, right=1cm, top=1.27cm, bottom=1.27cm]{geometry}
\usepackage{setspace}
\usepackage{tabularx}
\usepackage{makecell}
\usepackage{longtable}
\usepackage{multirow}
\usepackage{caption}
\usepackage{amsmath}
\usepackage{amssymb}
\usepackage{ragged2e}
\onehalfspacing


\begin{document}


\begin{table}[h]
    \centering
    \renewcommand{\arraystretch}{1.5}
    \begin{tabular}{m{11cm}m{6cm}}
        \begin{center}
            \textbf{Университет ИТМО}\\
            \textbf{Физико-технический мегафакультет}\\
            \textbf{Физический факультет}
        \end{center} 
        & 
        \begin{center}
            
\includegraphics[height=1.3cm]{logo.jpg}
        \end{center} \\
    \end{tabular}
\end{table}

\noindent\makebox[\linewidth]{\rule{\textwidth}{2pt}}

\vspace{10pt}

\begin{table}[h]
    \centering
    \renewcommand{\arraystretch}{1.7}
    \begin{tabular}{p{4cm}p{4cm}p{8cm}}
        Группа N3250 & & К работе допущен \underline{\hspace{4cm}} \\
        Студенты & \hfill \multirow[0pt]{3}{*}{\makecell[c]{Пахомов В.Ю \\ Булин М.М \\ Рулев М.Ю}}  & Работа выполнена \underline{\hspace{4cm}} \\
        & & \\
        & & \\
        \multicolumn{2}{l}{Преподаватель Рудель А.Е}  & Отчет принят \underline{\hspace{5cm}} \\
    \end{tabular}
\end{table}

\begin{center}
    \fontsize{20pt}{21.5pt} \selectfont
    \textbf{Рабочий протокол и отчет по} \\
    \textbf{лабораторной работе № 3.11}
\end{center}

\noindent\makebox[\linewidth]{\rule{\textwidth}{1pt}}
\noindent\makebox[\linewidth]{\rule{\textwidth}{1pt}}

\begin{justify}
    \textbf{1. Цель работы} \\
    1) Изучение вынужденных колебаний и явления резонанса напряжений в последовательном колебательном контуре. \\
    2) Построение резонансной кривой и определение резонансной частоты. \\
    3) Определение активного сопротивления и добротности колебательного контура. \\
\end{justify}

\begin{justify}
    \textbf{2. Задачи, решаемые при выполнении работы} \\
    1) Собрать электрическую цепь, состоящую из последовательно соединенных генератора гармонических колебаний, катушки индуктивности, магазина емкостей и активного сопротивления. \\
    2) Экспериментально снять амплитудно-частотную характеристику (АЧХ) контура: зависимость напряжения на конденсаторе от частоты входного сигнала для емкости $C=0{,}1$ мкФ. \\
    3) Определить по графику резонансную частоту и ширину полосы пропускания, рассчитать добротность контура разными методами. \\
    4) Провести серию измерений резонансных частот для различных значений емкости (от 1 нФ до 300 нФ). \\
    5) Построить линеаризованный график зависимости $f^2(1/C)$ и с его помощью определить индуктивность катушки $L$ и активное сопротивление контура $R$.
\end{justify}

\begin{justify}
    \textbf{3. Объект исследования} \\
    Объектом исследования является последовательный колебательный контур ($RLC$-контур), состоящий из катушки индуктивности $L$, конденсатора переменной емкости $C$ и активного сопротивления $R$, подключенный к источнику внешней гармонической ЭДС.
\end{justify}

\begin{justify}
    \textbf{4. Метод экспериментального исследования} \\
    Метод заключается в изучении отклика колебательной системы на внешнее периодическое воздействие. Изменяя частоту вынуждающей ЭДС (синусоидального сигнала с генератора) при постоянной амплитуде, фиксируется изменение амплитуды напряжения на емкости. Явление резонанса напряжений регистрируется по максимуму показаний осциллографа. Параметры контура ($L, R, Q$) определяются косвенным путем на основе анализа формы резонансной кривой и зависимости резонансной частоты от параметров элементов цепи.
\end{justify}

\begin{justify}
    \textbf{5. Рабочие формулы и исходные данные} \\
    $C_1 = 0.022 \; \text{мкФ} $
    $C_2 = 0.033 \; \text{мкФ} $
    $C_3 = 0.047 \; \text{мкФ} $
    $C_4 = 0.47 \; \text{мкФ} \pm 10\% $ \\
    $S = 64 \pm 0.05 \; \text{см} $
    $L = 7.8 \pm 0.1 \; \text{см} $
    $N1 = 1700 \; \text{вит} $
    $N2 = 1000 \; \text{вит} $ \\
    $R_1 = 65 \; \text{Ом} $
    $R_2 = 150 \; \text{Ом} $
    $R_3 = 680 \; \text{Ом} $
    $R_4 = 820 \; \text{Ом} $
    $R_5 = 470 \; \text{кОм} \pm 10\% $
    $L = 10 \; \text{мГн} \pm 10\% $ \\
    1. Формула Томсона (собственная частота контура без затухания):
    \[ f_0 = \frac{1}{2\pi\sqrt{LC}} \]
    2. Резонансная частота напряжения на конденсаторе (с учетом затухания):
    \[ f_{\text{рез}} = \sqrt{f_0^2 - 2\left(\frac{R}{4\pi L}\right)^2} \]
    3. Уравнение для графического определения индуктивности (из методических указаний):
    \[ f_{\text{рез}}^2 = \frac{1}{4\pi^2 L} \cdot \frac{1}{C} - \frac{R^2}{16\pi^2 L^2} \]
    4. Добротность контура (через отношение напряжений):
    \[ Q = \frac{U_{C,\text{рез}}}{\mathcal{E}_0} \]
    5. Добротность контура (через ширину полосы пропускания):
    \[ Q = \frac{f_{\text{рез}}}{\Delta f}, \quad \text{где } \Delta f = f_2 - f_1 \text{ на уровне } \frac{U_{\text{макс}}}{\sqrt{2}} \]
    6. Расчет активного сопротивления контура:
    \[ R = \frac{2\pi f_{\text{рез}} L}{Q} \]
\end{justify}


\begin{flushleft}
    \textbf{6. Измерительные приборы} \\
    \begin{table}[h]
    \centering
    \renewcommand{\arraystretch}{1.5}
    \begin{tabular}{|c|c|c|c|c|}
        \hline
        № п/п & Наименование & Тип прибора & Используемый диапазон & Погрешность прибора\\
        \hline
        1 & С3-ЭМ01 & Стенд & 0-20 В & 0.1 В \\
        \hline
        2 & АКИП-3409/2 & Генератор сигналов & 20-40 Гц & 0.1 Гц \\
        \hline
        3 & ФПЭ-11 & Осцилограф & 0-20 В & 0.1В \\
        \hline
    \end{tabular}
    \end{table}
\end{flushleft}


\begin{flushleft}
    \textbf{7. Схема установки} \\
    \begin{figure}[h!]
        \centering
        \includegraphics[height=8cm]{schema1.png}
        \caption{Общий вид магнитной установки}
    \end{figure}
    \begin{figure}[h!]
        \centering
        \includegraphics[height=10cm]{schema2.png}
        \caption{Общая схема магнитной установки}
    \end{figure}
\end{flushleft}


\newpage


\begin{flushleft}
    \textbf{8. Результаты прямых измерений и их обработки} \\
    \begin{table}[h]
    \centering
    \caption{Экспериментальные данные амплитудно-частотной характеристики}
    \vspace{0.3cm}
    \begin{tabular}{|c|c||c|c|}
    \hline
    \textbf{Частота, Гц} & \textbf{Амплитуда, В} & \textbf{Частота, Гц} & \textbf{Амплитуда, В} \\ \hline
    4000 & 7,60 & 5100 & 8,24 \\ \hline
    4100 & 7,76 & 5200 & 7,92 \\ \hline
    4200 & 8,08 & 5300 & 7,76 \\ \hline
    4300 & 8,24 & 5400 & 7,44 \\ \hline
    4400 & 8,40 & 5500 & 7,12 \\ \hline
    4500 & 8,56 & 5600 & 6,00 \\ \hline
    4600 & 8,56 & 5700 & 6,48 \\ \hline
    \textbf{4700} & \textbf{8,72} & 5800 & 6,16 \\ \hline
    4800 & 8,56 & 5900 & 5,76 \\ \hline
    4900 & 8,56 & 6000 & 5,52 \\ \hline
    5000 & 8,40 &      &      \\ \hline
    \end{tabular}
    \end{table}

    
\end{flushleft}


\newpage


\begin{justify}
    \textbf{9. Расчет результатов косвенных измерений (таблицы, примеры расчетов).} \\

    Из таблицы и графика видно, что экспериментальная резонансная частота:
    \[ f_{\text{эксп}} = 4700 \text{ Гц} \]

    Максимальная амплитуда напряжения на конденсаторе:
    \[ U_{\text{макс}} = 8{,}72 \text{ В} \]

    Расчетная резонансная частота определяется по формуле Томсона:
    \[ f_{\text{расч}} = \frac{1}{2\pi\sqrt{LC}} = \frac{1}{2\pi\sqrt{10 \times 10^{-3} \times 0{,}1 \times 10^{-6}}} = \frac{1}{2\pi \times 1 \times 10^{-4}} = 5032{,}9 \text{ Гц} \]

    Экспериментальная резонансная частота (4700 Гц) оказалась ниже расчетной (5033 Гц). Это расхождение объясняется влиянием активного сопротивления контура.

    \[ \Omega_{C,\text{рез}} = \sqrt{\omega_0^2 - 2\beta^2} = \omega_0\sqrt{1 - 2\left(\frac{\beta}{\omega_0}\right)^2} \]
    где $\beta = R/(2L)$ — коэффициент затухания.

    При этом резонансная частота смещается в сторону меньших значений по сравнению с собственной частотой $\omega_0$. Чем больше активное сопротивление контура, тем больше смещение резонансной частоты и тем более размытым становится резонансный пик. \\
    Добротность по ширине резонансной кривой:
    \[ Q = \frac{f_{\text{рез}}}{\Delta f} \]
    где $\Delta f$ — ширина резонансной кривой на уровне 0,707 от максимума (уровень -3 дБ).

    Уровень -3 дБ соответствует амплитуде:
    \[ U_{-3\text{дБ}} = \frac{U_{\text{макс}}}{\sqrt{2}} = \frac{8{,}72}{1{,}414} = 6{,}166 \text{ В} \]

    Для определения ширины полосы пропускания $\Delta f$ необходимо найти частоты $f_1$ и $f_2$, при которых амплитуда равна 6,166 В. \\
    Слева от резонанса:
    \begin{itemize}
        \item При $f = 4000 \text{ Гц}$: $U = 7{,}60 \text{ В} > 6{,}166 \text{ В}$
        \item Точка входа в полосу: приблизительно $f_1 \approx 3900 \text{ Гц}$
    \end{itemize}

    Справа от резонанса:
    \begin{itemize}
        \item При $f = 5800 \text{ Гц}$: $U = 6{,}16 \text{ В} \approx 6{,}166 \text{ В}$
        \item Точка выхода из полосы: $f_2 \approx 5800 \text{ Гц}$
    \end{itemize}
    \[ \Delta f = f_2 - f_1 \approx 5800 - 3900 = 1900 \text{ Гц} \]
    \[ Q = \frac{f_{\text{рез}}}{\Delta f} = \frac{4700}{1900} = 2{,}47 \]
    Доброность контура по напряжению: \\
    $U_{C,\text{рез}} = 8{,}72 \text{ В}$ \\
    $E_0 = 4{,}0 \text{ В}$ \\
    \[ Q = \frac{U_{C,\text{рез}}}{E_0} = \frac{8{,}72}{4{,}0} = 2{,}18 \]
    Оба значения добротностей находятся в пределах 10-20\% друг от друга, что свидетельствует о хорошей согласованности результатов.\\
    Для построения графика зависимости квадрата резонансной частоты от обратной емкости воспользуемся формулой:
    \[ f_{\text{рез}}^2 = \frac{1}{4\pi^2 L} \cdot \frac{1}{C} - \frac{R^2}{16\pi^2 L^2} \]
    Это уравнение прямой вида $f^2 = k \cdot (1/C) + b$, где угловой коэффициент $k = \frac{1}{4\pi^2 L}$.
    \begin{table}[!h]
    \centering
    \small
    \begin{tabular}{|l|c|c|c|c|}
    \hline
    \textbf{Емкость} & \textbf{C, Ф} & \textbf{$f_{\text{теор}}$, Гц} & \textbf{$f_{\text{эксп}}$, Гц} & \textbf{$U_{\text{макс}}$, В} \\ \hline
    1 нФ & $1{,}0\times10^{-9}$ & 50329,2 & 57000 & 17,10 \\ \hline
    3 нФ & $3{,}0\times10^{-9}$ & 29057,6 & 34000 & 17,50 \\ \hline
    10 нФ & $1{,}0\times10^{-8}$ & 15915,5 & 18000 & 17,20 \\ \hline
    30 нФ & $3{,}0\times10^{-8}$ & 9188,8 & 9000 & 14,20 \\ \hline
    100 нФ & $1{,}0\times10^{-7}$ & 5032,9 & 4700 & 8,72 \\ \hline
    300 нФ & $3{,}0\times10^{-7}$ & 2905,8 & 2000 & 5,04 \\ \hline
    \end{tabular}
    \end{table}

    По методу наименьших квадратов получен угловой коэффициент:
    \[ k \approx 2{,}533 \times 10^6 \;\text{Гц}^2 \cdot \text{Ф} \]

    Из углового коэффициента находим индуктивность:
    \[ L = \frac{1}{4\pi^2 k} = \frac{1}{4\pi^2 \times 2{,}533 \times 10^6} \approx 0{,}01 \text{ Гн} = 10{,}00 \text{ мГн} \]
    Экспериментально определенная индуктивность совпадает с заданным значением.

    Оценка активного сопротивления контура:
    \[ R = \frac{2\pi f_{\text{рез}} L}{Q} =  \frac{2\pi \times 4700 \times 0{,}01}{2{,}18} \approx 135{,}5 \text{ Ом} \]
    Полученное значение $R \approx 135{,}5 \; \text{Ом}$ близко к номиналу включенного в цепь резистора $R_2 = 150 \; \text{Ом}$.

\end{justify}


\newpage


\begin{justify}
    \textbf{10. Расчет погрешностей измерений (для прямых и косвенных измерений)}\\
    Относительная погрешность добротности:
    \[ \varepsilon_Q = \sqrt{\varepsilon_{U_C}^2 + \varepsilon_{E_0}^2} = \sqrt{0{,}05^2 + 0{,}05^2} \approx 0{,}071 \; (7{,}1\%) \]
    Абсолютная погрешность:
    \[ \Delta Q = Q \cdot \varepsilon_Q = 2{,}18 \cdot 0{,}071 \approx 0{,}15 \]
    Результат: $Q = 2{,}2 \pm 0{,}2$.
    \[ \varepsilon_f = \frac{\Delta f}{f_{\text{рез}}} = \frac{100}{4700} \approx 0{,}021 \; (2{,}1\%) \]
    Погрешность индуктивности $L$:
    \[ \Delta L = L \cdot \varepsilon_L = 10{,}00 \cdot 0{,}10 = 1{,}00 \text{ мГн} \]
    Результат: $L = 10{,}0 \pm 1{,}0 \text{ мГн}$. \\
    Погрешность активного сопротивления $R$: \\
    \[ \varepsilon_R = \sqrt{\varepsilon_f^2 + \varepsilon_L^2 + \varepsilon_Q^2} \]
    \[ \varepsilon_R = \sqrt{0{,}021^2 + 0{,}10^2 + 0{,}071^2} = \sqrt{0{,}0004 + 0{,}01 + 0{,}005} \approx \sqrt{0{,}0154} \approx 0{,}124 \; (12{,}4\%) \]
    Абсолютная погрешность для $R \approx 135{,}5$ Ом:
    \[ \Delta R = R \cdot \varepsilon_R = 135{,}5 \cdot 0{,}124 \approx 16{,}8 \text{ Ом} \]
    Результат: $R = 136 \pm 17 \text{ Ом}$.
    \[ \delta f = \frac{|f_{\text{расч}} - f_{\text{эксп}}|}{f_{\text{расч}}} \times 100\% = \frac{|5032{,}9 - 4700|}{5032{,}9} \times 100\% \approx 6{,}6\% \]
\end{justify}


\newpage


\begin{justify}
    \textbf{11. Графики (перечень графиков, которые составляют Приложение 2)} \\
    \begin{figure}[h!]
        \centering
        \includegraphics[height=10cm]{graphic3.png} \\
        \caption{График зависимости резонансной кривой напряжения на конденсаторе }
    \end{figure}
    
    \begin{figure}[h!]
    \centering
    \includegraphics[height=10cm]{graphic2.png} \\
    \caption{График зависимости $\Omega_{res}^2$ ($\text{рад}^2/c^2$) }
    \end{figure}
\end{justify}


\newpage


\begin{justify}
    \textbf{12. Окончательные результаты} \\
    \begin{table}[h]
    \centering
    \begin{tabular}{|c|l|c|c|}
    \hline
    \textbf{№} & \textbf{Параметр} & \textbf{Значение} & \textbf{Ед. изм.} \\ \hline
    1 & Расчётная резонансная частота & 5032,9 & Гц \\ \hline
    2 & Экспериментная резонансная частота & 4700 & Гц \\ \hline
    3 & Максимальное напряжение на $C$ & 8,72 & В \\ \hline
    4 & Входное напряжение & 4,0 & В \\ \hline
    5 & Добротность (по напряжению) & 2,18 & — \\ \hline
    6 & Добротность (по ширине) & 2,47 & — \\ \hline
    7 & Полоса пропускания -3 дБ & 1900 & Гц \\ \hline
    8 & Заданная индуктивность & 10,0 & мГн \\ \hline
    9 & Экспериментальная индуктивность & 10,00 & мГн \\ \hline
    10 & Активное сопротивление & 135 & Ом \\ \hline
    11 & Коэффициент затухания & 6,75 & с$^{-1}$ \\ \hline
    \end{tabular}
    \end{table}
\end{justify}

\begin{justify}
    \textbf{13. Выводы и анализ результатов работы}

    В ходе выполнения лабораторной работы были изучены вынужденные электромагнитные колебания в последовательном $RLC$-контуре. На основе полученных экспериментальных данных сделаны следующие выводы:

    \begin{enumerate}
        \item Исследование резонанса напряжений: \\
        Экспериментально получена амплитудно-частотная характеристика контура. Резонанс напряжений на конденсаторе наблюдался на частоте $f_{\text{эксп}} = 4700$ Гц. Это значение оказалось ниже теоретической собственной частоты контура $f_0 \approx 5033$ Гц (отклонение 6,6\%). Наблюдаемое смещение резонансного максимума в область низких частот (<<влево>>) полностью согласуется с теорией для контуров с заметным затуханием ($\Omega_{\text{рез}} < \omega_0$).

        \item Определение добротности: \\
        Добротность контура была оценена двумя независимыми методами:
        \begin{itemize}
            \item По крастности перенапряжения: $Q_U = 2{,}18 \pm 0{,}15$.
            \item По ширине полосы пропускания: $Q_{\Delta f} \approx 2{,}47$.
        \end{itemize}
        Значения совпадают в пределах погрешности измерений. Низкая величина добротности ($Q \sim 2$) свидетельствует о значительных активных потерях в цепи.

        \item Определение индуктивности: \\
        Построен график зависимости $f_{\text{рез}}^2$ от обратной емкости $1/C$. Линейный характер зависимости подтвердил справедливость формулы Томсона. Рассчитанная по угловому коэффициенту графика индуктивность составила $L = 10{,}0 \pm 1{,}0$ мГн, что совпадает с номиналом используемой катушки.

        \item Определение активного сопротивления: \\
        Рассчитанное полное активное сопротивление контура составило $R = 136 \pm 17$ Ом. Это значение с высокой точностью соответствует номиналу резистора $R_2 = 150$ Ом (допуск $\pm 5\%$), который был включен в цепь. Незначительное расхождение объясняется влиянием внутреннего сопротивления генератора и погрешностями метода определения добротности.
    \end{enumerate}

\end{justify}

\begin{justify}
    \textbf{14. Дополнительные задания}
\end{justify}

\begin{justify}
    \textbf{15. Выполнение дополнительных заданий}
\end{justify}

\begin{justify}
    \textbf{16. Замечания преподавателя (исправления, вызванные замечаниями преподавателя, также помещают в этот пункт)}
\end{justify}

\end{document}
